\documentclass{article}
\usepackage[utf8]{inputenc}
\usepackage{geometry}
\geometry{a4paper, margin=1in}
\usepackage{parskip}
\usepackage{hyperref}

\title{The Collision of Freedoms: A Structural Problem}
\author{Albert Jan van Hoek}
\date{\today}

\begin{document}

\maketitle

\section*{The Collision of Freedoms}

Two types of freedom are colliding right now — and most of us are on the losing side.

The first is the freedom we usually mean: \textbf{doing whatever you (singular) want.}

The second is rarer, harder, and more powerful: \textbf{doing whatever you (plural) decide together.}

In my \textbf{Theory of Long-term Collaboration (TLC)}, this collision is the central problem of any shared system — from a marriage to a civilization.

Collaboration is simply more cumbersome than domination.

\begin{itemize}
    \item Listening takes time.
    \item Adapting takes energy.
    \item Listening to women takes time.
    \item Listening to minorities takes time.
    \item Listening to what the planet needs takes time.
\end{itemize}

Anyone who has ever had to truly share a decision knows this in their bones.

And right now, we are all experiencing what it means when that sharing breaks down.

The decisions that shape our collective future are being made by a handful of people — people like you and me, bags of bones and blood, nothing mystical about them — who have discovered they can simply ignore the second freedom entirely.

Do whatever they want. Leave the rest of us to beg, shout, or hope for sympathy.

\section*{A Universal Problem}

This is not a new problem. It appears at every scale.
Even in a marriage, the struggle over whose perspective counts — whose freedom gets to define the shared future — can quietly become a de facto dictatorship.

The scale changes. The structure doesn't.

\section*{What Does TLC Say to Do?}

Not revolution. Not surrender.

Something harder: \textbf{change the structure itself, without breaking it, without force.}
Make the second freedom — collective decision-making — too costly to ignore.

That starts with seeing the problem clearly for what it is: not a values failure, not a lack of sympathy, but a \textbf{structural imbalance between two kinds of freedom.}

Only one of them scales. The other eventually becomes domination.

\end{document}
